\chapter{Desarrollo}

\section{Entorno de desarrollo}
% [PENDIENTE: detallar versiones: Python 3.11, OpenCV 4.x, Node.js, etc.]

\section{Implementaci\'{o}n del m\'{o}dulo de seguimiento (tracker.py)}

\subsection{Representación de imágenes y video}

Para el módulo de seguimiento, un video se trata como una secuencia de matrices de
números. Cada cuadro (\textit{frame}) es una imagen BGR de tres canales (azul, verde,
rojo), que es la convención de OpenCV. Un video de 5 minutos a 30 FPS equivale a
procesar 9{,}000 cuadros de forma secuencial.

\subsection{Modelo de fondo y sustracción}

El primer paso del pipeline es construir un modelo del fondo: una imagen de referencia
del cilindro sin la rata. Se calcula como la mediana píxel a píxel de los primeros
60 cuadros del video. Usar la mediana en vez del promedio evita que una rata inmóvil
al inicio ``contamine'' el modelo. Después, en cada cuadro se calcula la diferencia
absoluta contra esa referencia; los píxeles que cambiaron significativamente son el
primer plano (la rata).

Este enfoque es más simple que métodos adaptativos como el modelo de mezcla gaussiana
(GMM) \cite{zivkovic2004}, pero es suficiente para este caso de uso: la cámara es fija
y el fondo no cambia durante el experimento.

\subsection{Umbralización y limpieza morfológica}

La imagen de diferencias se convierte a binario con el método de Otsu \cite{otsu1979},
que calcula automáticamente el umbral óptimo a partir del histograma de intensidades.
Después se aplica una apertura morfológica (erosión seguida de dilatación) para
eliminar ruido pequeño y rellenar huecos en la región detectada.

\subsection{Detección de contornos y selección de candidatos}

Sobre la imagen binaria limpia se extraen los contornos con \texttt{findContours} de
OpenCV \cite{suzuki1985}. De cada contorno se obtiene el área, el rectángulo envolvente
(\textit{bounding box}) y el centroide. Con esas propiedades se filtra cuál de los
candidatos detectados corresponde a la rata y cuál es ruido, usando restricciones de
área mínima y máxima calibradas para el tamaño esperado del animal dentro de cada
cilindro.

Para verificar consistencia entre cuadros consecutivos se usa la métrica IoU
(\textit{Intersection over Union}): la relación entre el área de intersección y el
área de unión de dos rectángulos. Un IoU cercano a 1 indica que la detección actual
coincide con la del cuadro anterior; si baja demasiado se activa el respaldo del
tracker de correlación.

\subsection{Estimación de la línea de agua}

La superficie del agua puede generar reflejos que hacen que el animal aparezca
duplicado desde la vista lateral. Para mitigar ese efecto, el sistema estima en cada
cuadro la posición vertical de la línea del agua: analiza el perfil de intensidad
promedio por filas dentro de cada ROI e identifica la fila con el cambio más brusco.
Como esa estimación oscila por ruido entre cuadros, se estabiliza con un promedio
móvil exponencial (EMA):

\begin{equation}
    \hat{y}_t = \alpha \cdot \tilde{y}_t + (1 - \alpha) \cdot \hat{y}_{t-1}
    \label{eq:ema}
\end{equation}

donde $\tilde{y}_t$ es la estimación del cuadro actual, $\hat{y}_{t-1}$ es el valor
suavizado anterior, y $\alpha$ controla la velocidad de respuesta. Esta lógica está
encapsulada en la clase \texttt{WaterlineSmoother}. Su efectividad depende de las
condiciones del video y se evaluará durante las pruebas de validación.

% [PENDIENTE: describir las clases principales: WaterlineSmoother, RatTracker]
% [PENDIENTE: fragmentos de c\'{o}digo relevantes y su explicaci\'{o}n]

\section{Implementaci\'{o}n del clasificador de conducta}
% [PENDIENTE: pendiente Sprint 4 --- arquitectura del clasificador,
%   proceso de fine-tuning, conjunto de entrenamiento y anotaci\'{o}n con BORIS]

\section{Implementaci\'{o}n del backend (Flask)}
% [PENDIENTE: endpoints implementados, modelo de datos SQLAlchemy]

\section{Implementaci\'{o}n del frontend (React + Vite)}
% [PENDIENTE: vistas implementadas, comunicaci\'{o}n con API]

\section{Integraci\'{o}n y despliegue con Docker Compose}
% [PENDIENTE: descripci\'{o}n del docker-compose.yml y servicios definidos]